% GENERATED FILE. Edit canonical lessons, then run npm run generate:books.
% canonical-source-hash: fnv1a64:94caafaa94697a72
% canonical-lessons: SA-C42-sometimes, SA-C42-suddenly, SA-C42-mostly, SA-C42-a-little, SA-C42-much

\chapter{Five More Replies}
\label{ch:sa-replies-third}
\hlchaptermodality{42}

\begin{chapteropening}
\textbf{By the end of this chapter:} \emph{I can say sometimes, suddenly, mostly, a little and much in Sanskrit, and use any of the five on its own as a reply.}

Complete SA-C42-much and demonstrate the chapter promise.
\end{chapteropening}

\section[kadācit]{\sk{कदाचित्} (kadācit) — sometimes}

\label{lesson:SA-C42-sometimes}



\begin{quote}
Before the new word: what did \emph{chatram} mean?
\end{quote}

\subsection*{You'll want to know: \sk{कदाचित्}}
\textbf{\sk{कदाचित्}} (\emph{kadācit}) — \textquotedblleft{}sometimes\textquotedblright{}.

\textbf{\sk{कदाचित्}} answers \textquotedblleft{}when?\textquotedblright{} with \textquotedblleft{}now and then\textquotedblright{}.

Both pieces are yours already: \sk{कदा}, the \emph{when} you ask questions with, and a small \emph{cit} fastened on the end. \textbf{\sk{कदाचित्}} is \emph{when} gone vague — at some time or other, and no promise which.

It is also the polite way to leave a door open: not yes, not no, perhaps sometimes.

\begin{grammarlens}[title={What you've built}]
A reply that answers \emph{when} without saying when.
\end{grammarlens}

\subsection*{Your turn}
Say these aloud:

\begin{itemize}
\raggedright
  \item \emph{kadācit}
  \item it once more, slowly
  \item \emph{chatram}, then \emph{kadācit}
\end{itemize}

\begin{tcolorbox}[breakable,colback=teal!4,colframe=teal!35!black,title={Before you move on}]
What does \sk{कदाचित्} mean? (\textquotedblleft{}sometimes\textquotedblright{}.) Say it once more.
\end{tcolorbox}
\section[sahasā]{\sk{सहसा} (sahasā) — suddenly}

\label{lesson:SA-C42-suddenly}



\begin{quote}
Before the new word: what did \emph{kadācit} mean?
\end{quote}

\subsection*{You'll want to know: \sk{सहसा}}
\textbf{\sk{सहसा}} (\emph{sahasā}) — \textquotedblleft{}suddenly\textquotedblright{}.

\textbf{\sk{सहसा}} is \textquotedblleft{}all at once, without warning\textquotedblright{}.

Underneath it is \emph{sahas}, \textquotedblleft{}force, a sudden strength\textquotedblright{}, put into the shape that means \emph{by means of}: by force, and so at a rush. \sk{कृपया}, the please you already say, is made the very same way out of \emph{kṛpā}.

It carries a shade of judgement with it — a thing done \textbf{\sk{सहसा}} is a thing done without thinking first.

\begin{grammarlens}[title={What you've built}]
One thing happening at a rush, and a mild opinion about it.
\end{grammarlens}

\subsection*{Your turn}
Say these aloud:

\begin{itemize}
\raggedright
  \item \emph{sahasā}
  \item it once more, slowly
  \item \emph{kadācit}, then \emph{sahasā}
\end{itemize}

\begin{tcolorbox}[breakable,colback=teal!4,colframe=teal!35!black,title={Before you move on}]
What does \sk{सहसा} mean? (\textquotedblleft{}suddenly\textquotedblright{}.) Say it once more.
\end{tcolorbox}
\section[prāyaḥ]{\sk{प्रायः} (prāyaḥ) — mostly, for the most part}

\label{lesson:SA-C42-mostly}



\begin{quote}
Before the new word: what did \emph{sahasā} mean?
\end{quote}

\subsection*{You'll want to know: \sk{प्रायः}}
\textbf{\sk{प्रायः}} (\emph{prāyaḥ}) — \textquotedblleft{}mostly, for the most part\textquotedblright{}.

\textbf{\sk{प्रायः}} is \textquotedblleft{}for the most part, in the main, usually\textquotedblright{}.

Say it when something holds most of the time and you would rather not promise more than that. It is \emph{pra-}, \textquotedblleft{}forward\textquotedblright{}, with the going-piece you met in \sk{विहगः}: what runs out in front, what happens in the main.

Beside \sk{सर्वम्}, which claims everything, \textbf{\sk{प्रायः}} claims almost everything and keeps its hedge.

\begin{grammarlens}[title={What you've built}]
A reply that is honest about not being certain.
\end{grammarlens}

\subsection*{Your turn}
Say these aloud:

\begin{itemize}
\raggedright
  \item \emph{prāyaḥ}
  \item it once more, slowly
  \item \emph{sahasā}, then \emph{prāyaḥ}
\end{itemize}

\begin{tcolorbox}[breakable,colback=teal!4,colframe=teal!35!black,title={Before you move on}]
What does \sk{प्रायः} mean? (\textquotedblleft{}mostly, for the most part\textquotedblright{}.) Say it once more.
\end{tcolorbox}
\section[kiñcit]{\sk{किञ्चित्} (kiñcit) — a little, somewhat}

\label{lesson:SA-C42-a-little}



\begin{quote}
Before the new word: what did \emph{prāyaḥ} mean?
\end{quote}

\subsection*{You'll want to know: \sk{किञ्चित्}}
\textbf{\sk{किञ्चित्}} (\emph{kiñcit}) — \textquotedblleft{}a little, somewhat\textquotedblright{}.

\textbf{\sk{किञ्चित्}} is \textquotedblleft{}a bit, somewhat, a small amount\textquotedblright{}.

It is \sk{किम्}, the \emph{what} you have been asking with since your first questions, wearing the same little \emph{cit} that softened \textbf{\sk{कदाचित्}}. \emph{What} gone vague is \emph{something} — and \emph{something} said of a quantity is \emph{a little}.

\sk{अल्पम्} describes a small thing; \textbf{\sk{किञ्चित्}} measures a small amount of it.

\begin{grammarlens}[title={What you've built}]
A small helping of an answer.
\end{grammarlens}

\subsection*{Your turn}
Say these aloud:

\begin{itemize}
\raggedright
  \item \emph{kiñcit}
  \item it once more, slowly
  \item \emph{prāyaḥ}, then \emph{kiñcit}
\end{itemize}

\begin{tcolorbox}[breakable,colback=teal!4,colframe=teal!35!black,title={Before you move on}]
What does \sk{किञ्चित्} mean? (\textquotedblleft{}a little, somewhat\textquotedblright{}.) Say it once more.
\end{tcolorbox}
\section[bahu]{\sk{बहु} (bahu) — much, many}

\label{lesson:SA-C42-much}



\begin{quote}
Before the last word of the chapter: what did \emph{prāyaḥ} mean, and what did \emph{kiñcit} mean?
\end{quote}

\subsection*{You'll want to know: \sk{बहु}}
\textbf{\sk{बहु}} (\emph{bahu}) — \textquotedblleft{}much, many\textquotedblright{}.

\textbf{\sk{बहु}} is \textquotedblleft{}much\textquotedblright{} of one thing and \textquotedblleft{}many\textquotedblright{} of several — the far end of the answer \textbf{\sk{किञ्चित्}} starts.

Greek \emph{pakhus}, \textquotedblleft{}thick, stout\textquotedblright{}, is its cousin, and thickness is the sense underneath: plenty as bulk rather than as a count.

Hindi \emph{bahut}, \textquotedblleft{}very, a lot\textquotedblright{}, is \textbf{\sk{बहु}} with an ending on it, and it is one of the commonest words anyone says today.

\begin{grammarlens}[title={What you've built}]
Five replies about how much and how often: sometimes, suddenly, mostly, a little, much.
\end{grammarlens}

\subsection*{Your turn}
Say these aloud:

\begin{itemize}
\raggedright
  \item \emph{bahu}
  \item it once more, slowly
  \item all five in order, then \emph{kadācit} and \emph{bahu} together
\end{itemize}

\begin{tcolorbox}[breakable,colback=teal!4,colframe=teal!35!black,title={Before you move on}]
What does \sk{बहु} mean? (\textquotedblleft{}much, many\textquotedblright{}.) Say it once more.
\end{tcolorbox}
